% ============================================================
% Appendix: Edge-by-Edge Evidence Table for the Causal DAG
% Project: Energy Efficiency and Mental Health at the LSOA Level
% ============================================================
% Usage: \input{appendix_table.tex} within your main document
% Requires: longtable, booktabs, array, caption, makecell packages
% ============================================================

\section{Edge-by-Edge Evidence for the Causal DAG}\label{app:edge-table}

Tables~\ref{tab:edges-hee}--\ref{tab:edges-social} present the full evidence base for each directed edge in the DAG. For each edge, the table provides a brief causal rationale and one or more supporting references. Edges are grouped thematically to align with the narrative review in Section~\ref{sec:dag-literature}.

% =====================================================
% Edges into HEE (Exposure)
% =====================================================

\begin{longtable}{@{}p{3.5cm}p{11.5cm}@{}}
\caption{Edges into Home Energy Efficiency (HEE).}
\label{tab:edges-hee}\\
\toprule
\textbf{Edge} & \textbf{Rationale and supporting evidence} \\
\midrule
\endfirsthead
\multicolumn{2}{l}{\textit{Table \ref{tab:edges-hee} continued from previous page}}\\
\toprule
\textbf{Edge} & \textbf{Rationale and supporting evidence} \\
\midrule
\endhead
\bottomrule
\endfoot
Economic\_Capital $\to$ HEE & Household income constrains the capacity to invest in retrofit measures, which typically require substantial upfront capital even where government subsidies are available. Evidence from a controlled before-and-after study in low-income Welsh neighbourhoods demonstrates that energy efficiency investments are concentrated in areas targeted by income-based programmes \citep{grey_short-term_2017}. Quasi-experimental analysis of low-income English households shows that modest income increases translate directly into higher energy spending, consistent with a binding budget constraint on household energy behaviour \citep{galvin_juggling_2024}. \\
\addlinespace
Urbanicity $\to$ HEE & Rural dwellings in England have systematically lower EPC ratings than urban dwellings (median 59 versus 68), a gap only partially accounted for by housing stock composition. Analysis of 14~million English EPC records identifies rural-urban status as the single most important geographic predictor of dwelling-level energy efficiency, driven by older detached stock, absence of mains gas connection, and remoteness from retrofit services \citep{sheridan_energy_2023, office_for_national_statistics_energy_2025}. Modern high-density urban construction complying with updated building regulations further widens this gap for new-build stock. \\
\addlinespace
Policy\_Scheme\_Access $\to$ HEE & Household- and area-level uptake of government energy efficiency programmes directly determines whether retrofit measures are installed. The Energy Company Obligation (ECO, 2013--present, currently in its fourth iteration ECO4 running 2022--2026) has accounted for between 71\% and 86\% of measures installed under government schemes in each year since 2013, and for approximately 75\% of measures installed in 2024 \citep{department_for_energy_security_and_net_zero_household_2025}. The now-closed Green Deal (2013--2015) \citep{rosenow_post_2016} and the ongoing Warm Home Discount scheme further contribute to the distribution of HEE through means-tested and property-targeted eligibility criteria \citep{penasco_policy_2024}. \\
\addlinespace
Housing\_Characteristics $\to$ HEE & Physical dwelling form (age, construction era, built form, and materials) is the strongest single determinant of baseline energy efficiency. Pre-1919 English properties average EPC 60, compared with 74 for post-1990 builds \citep{office_for_national_statistics_energy_2025}. Solid-walled properties, listed buildings, and detached houses with high surface-to-volume ratios face both higher baseline heat loss and greater technical constraints on retrofit feasibility \citep{trotta_factors_2018, symonds_linking_2024}. \\
\addlinespace
Tenure $\to$ HEE & The landlord-tenant split incentive suppresses retrofit investment in the private rented sector: landlords bear the capital cost of improvements, while tenants capture the resulting bill savings, creating a principal-agent problem that depresses efficiency investment. Empirical evidence using US Residential Energy Consumption Survey data documents landlord underinvestment across multiple categories of residential energy efficiency, including space-heating systems, water-heating, window thickness, insulation, and weatherisation \citep{melvin_split_2018}. These findings are corroborated in the residential rental market by \citet{petrov_landlord-tenant_2021} and \citet{gillingham_split_2012}. \\
\end{longtable}


% =====================================================
% Edges into MH (Outcome)
% =====================================================
\begin{longtable}{@{}p{3.5cm}p{11.5cm}@{}}
\caption{Edges into Mental Health (MH).}
\label{tab:edges-mh}\\
\toprule
\textbf{Edge} & \textbf{Rationale and supporting evidence} \\
\midrule
\endfirsthead
\multicolumn{2}{l}{\textit{Table \ref{tab:edges-mh} continued from previous page}}\\
\toprule
\textbf{Edge} & \textbf{Rationale and supporting evidence} \\
\midrule
\endhead
\bottomrule
\endfoot
Education $\to$ MH & Higher educational attainment is associated with improved coping, health literacy, and self-efficacy, and is consistently linked to better adult mental health outcomes. Sibling fixed-effects analyses of UK Biobank data find that higher educational attainment remains significantly associated with fewer depressive symptoms after controlling for unobserved family-level confounders and polygenic scores, though the strength of these associations is reduced by 40–70\% relative to standard OLS estimates once sibling controls are applied \citep{amin_re-examining_2023}. Population-level evidence places education within the framework of fundamental social causation of mental health \citep{allen_social_2014, alqassim_impact_2024}. \\
\addlinespace
Economic\_Capital $\to$ MH & Poverty, unemployment, and financial stress are established causal risk factors for common mental disorders. Using panel data from the UK Household Longitudinal Study, \citet{thomson_effects_2023} estimate that falling below the poverty line increases common mental disorder prevalence by 2.5 percentage points. This finding is corroborated by systematic review and meta-analytic evidence synthesising causal studies of poverty and mental health in low- and middle-income as well as high-income contexts \citep{ridley_poverty_2020, salem_impact_2025}. \\
\addlinespace
Tenure $\to$ MH & Housing tenure insecurity, particularly in the private rented sector, generates anxiety through fear of eviction, uncontrolled rent increases, and limited control over dwelling conditions. Longitudinal analysis in the UK and international contexts demonstrates that housing insecurity independently predicts psychological distress, poor sleep, and hypertension, beyond the effect of income or housing physical quality \citep{thomson_impact_2024, rahman_unlocking_2024, zehrung_investigating_2024}. \\
\addlinespace
Overcrowding $\to$ MH & Residential overcrowding generates psychological distress through reduced privacy, interpersonal conflict, constrained personal space, and reduced control over the domestic environment. Longitudinal analysis of Chilean household trajectories finds that sustained overcrowding increases the probability of poor mental health, with effects attenuated but not eliminated after accounting for income \citep{ruiz-tagle_household_2022}. Parallel evidence from Sweden and New Zealand confirms the association in migrant and general populations \citep{mangrio_crowded_2018, pierse_examining_2016}. \\
\addlinespace
Climate $\to$ MH & Extreme weather events — particularly flooding and heatwaves — directly cause psychological trauma. A UK-focused meta-analysis of post-extreme-event mental health outcomes reports 12-month prevalence of 30\% for PTSD, 21\% for depression, and 20\% for anxiety among those directly affected \citep{cruz_effect_2020}. This edge captures the direct mental health consequences of climate exposure, distinct from indirect pathways through housing degradation and IEQ. \\
\addlinespace
Air\_Pollution $\to$ MH & Chronic exposure to ambient PM$_{2.5}$ and NO$_2$ is associated with increased depression and anxiety risk, proposed to operate through neuroinflammation and oxidative stress. A systematic review and meta-analysis identifies a pooled association of 10--18\% higher depression risk per 10\,µg/m$^{3}$ PM$_{2.5}$ \citep{braithwaite_air_2019}, corroborated in an updated synthesis by \citet{borroni_air_2022}, while \citet{bhui_air_2023} provides a contemporary review of the evidence base and mechanistic hypotheses. \\
\addlinespace
Green\_Space $\to$ MH & Access to and quality of green and blue space is protective against depression and anxiety, with proposed mechanisms including attention restoration and chronic stress reduction. A ten-year longitudinal dynamic panel study of 2.3 million adults in Wales finds significant protective associations between ambient greenness, local green space access, and subsequent mental health \citep{sheridan_mayall_ambient_2024}. Twin analyses and a recent meta-analysis further support the association \citep{cohen-cline_access_2015, liu_green_2023, kruize_urban_2019}. \\
\addlinespace
Noise $\to$ MH & Chronic environmental noise exposure — primarily from road, rail, and aircraft traffic — is hypothesised to affect mental health through annoyance, sleep disturbance, and sustained physiological stress. Evidence quality is mixed: a meta-analysis of road traffic noise and depression found a pooled 4\% increase in depression odds per 10\,dB(A) L$_\text{den}$ and 12\% for anxiety, though both estimates had confidence intervals crossing zero and the synthesis itself graded the underlying evidence as ``very low'' quality \citep{dzhambov_road_2019}. A more recent meta-analysis finds statistically significant effects specifically for aircraft noise (12\% increase in depression risk per 10\,dB L$_\text{den}$) but weaker, non-significant effects for road and rail traffic sources \citep{hegewald_traffic_2020}. \citet{hahad_noise_2024} offers a contemporary review integrating mechanistic evidence. \\
\addlinespace
IEQ $\to$ MH & Indoor damp, mould, cold, and indoor air pollutants are associated with distress, perceived loss of control, and measurable increases in depression and anxiety. A UKHSA scoping review finds that 88\% of included studies report significant associations between dampness/mould exposure and psychological outcomes including depression, anxiety, and perceived loss of control over the home environment \citep{brooks_psychological_2025}. \citet{shenassa_dampness_2007} establishes the dampness–depression pathway operates both through mould-related physical illness and through perceived control; \citet{liddell_living_2015} provides a framework for the cold-housing/MH mechanism. \\
\addlinespace
Social\_Capital $\to$ MH & Area-level social capital — community trust, civic participation, and associational density — is inversely associated with common mental disorder prevalence through pathways of social support, collective efficacy, and sense of belonging. A systematic review of cognitive social capital and mental disorder finds consistent protective associations across multiple study designs and contexts \citep{de_silva_social_2005}. Updated evidence confirms this across recent UK and international studies \citep{magro-montanes_relationship_2024, rodgers_social_2019}. \\
\addlinespace
Isolation $\to$ MH & Social isolation and loneliness are robust, independent risk factors for depression, anxiety, and cognitive decline. An umbrella review of systematic reviews finds that loneliness more than doubles the odds of developing new-onset depression across longitudinal studies \citep{leigh-hunt_overview_2017}, with \citet{holtlunstad_social_2024} providing a contemporary synthesis of the evidence and mechanistic pathways. \\
\addlinespace
Accessibility $\to$ MH & Physical access to healthcare, employment, and social services affects mental health both directly (through increased stress and reduced ability to obtain treatment) and indirectly (through participation in economic and social life). Evidence from UK health service inequalities research and international accessibility studies supports the association \citep{lowther-payne_understanding_2023, mseke_impact_2024}. \\
\addlinespace
Digital\_Capital $\to$ MH & Digital capital — broadband access, device availability, and digital literacy — has a net protective effect on mental health at the population level, primarily through access to mental health resources, telehealth, and maintenance of social connection. Digital exclusion disproportionately affects people with severe mental illness and older adults in deprived communities, compounding existing mental health inequalities \citep{greer_digital_2019, spanakis_digital_2021}. Despite adverse effects of digital overuse, at the LSOA aggregate level relevant to the present DAG, the access-to-services channel dominates. \\
\addlinespace
Substance\_Abuse $\to$ MH & Substance use disorders are highly comorbid with mental health disorders. A meta-analysis of epidemiological surveys finds that drug use disorder is associated with 3.8-fold increased odds of major depression, and alcohol use disorder with approximately 2.4-fold increased odds \citep{lai_prevalence_2015}. \\
\addlinespace
PH $\to$ MH & Chronic physical illness increases depression risk through pathways of disability, pain, inflammation, and reduced quality of life. A mediation analysis using UK longitudinal data demonstrates that physical health substantially mediates the income-mental-health relationship \citep{ohrnberger_relationship_2017}; a comprehensive recent review synthesises the comorbidity patterns between major depressive disorder and physical disease \citep{berk_comorbidity_2023}. The population-level burden is described in \citet{prince_no_2007}. \\
\addlinespace
Neighbourhood\_Stress $\to$ MH & Chronic neighbourhood stressors — principally crime and perceived unsafety — generate fear, hypervigilance, and sustained psychological distress. A meta-analysis of 63 studies finds significant pooled associations between neighbourhood crime and depression, anxiety, and general distress \citep{baranyi_impact_2021}. Non-crime dimensions of neighbourhood stress (physical disorder, anti-social behaviour) are included as part of the construct theoretically, though the dominant empirical evidence base operationalises neighbourhood stress through crime exposure. \\
\addlinespace
Demographics $\to$ MH & Age, gender, ethnicity, and household composition shape population-level distributions of common mental disorder. Women have substantially higher CMD prevalence than men in England (approximately 24\% versus 15\% in recent APMS data); ethnic minorities face discrimination-related and migration-related mental health stressors \citep{nhs_england_digital_adult_2024, irizar_prevalence_2024}. Fundamental-cause framing places these patterns within the broader social determinants of mental health \citep{allen_social_2014}. \\
\end{longtable}


% =====================================================
% Education and Economic Capital outgoing edges
% =====================================================

\begin{longtable}{@{}p{3.8cm}p{11.2cm}@{}}
\caption{Outgoing edges from Education and Economic Capital.}
\label{tab:edges-econ_edu}\\
\toprule
\textbf{Edge} & \textbf{Rationale and supporting evidence} \\
\midrule
\endfirsthead
\multicolumn{2}{l}{\textit{Table \ref{tab:edges-econ_edu} continued from previous page}}\\
\toprule
\textbf{Edge} & \textbf{Rationale and supporting evidence} \\
\midrule
\endhead
\bottomrule
\endfoot
\multicolumn{2}{@{}l}{\textit{Education outgoing edges}}\\
\addlinespace
Education $\to$ Economic\_Capital & Educational attainment causally affects adult earnings through labour market returns. UK-based instrumental variable estimates exploiting the 1947, 1963, and 1972 compulsory schooling reforms provide a robust case for approximately a 6\% rate of return to education for men \citep{dolton_returning_2017}. Broader international reviews of the returns-to-schooling literature report returns in the range of 5--10\% per additional year of schooling across developed economies, with UK-specific estimates clustering at the lower end of this range once reform-based identification strategies are applied \citep{card_causal_1999, blau_american_1978}. This edge also serves as the pathway through which education indirectly affects home energy efficiency, replacing the direct \textit{Education $\to$ HEE} edge excluded in Section~\ref{sec:excluded-edges}. \\
\addlinespace
Education $\to$ PH & Education is associated with reduced all-cause morbidity and mortality through pathways including health literacy, health behaviours, occupational exposures, and access to healthcare. A global systematic review and meta-analysis of 603 studies reports a reduction of approximately 1.9\% in adult mortality risk per additional year of schooling, with effects concentrated in causes of death amenable to health-behavioural mediation \citep{balaj_effects_2024}. \\
\addlinespace
\multicolumn{2}{@{}l}{\textit{Economic Capital outgoing edges}}\\
\addlinespace
Economic\_Capital $\to$ PH & Socioeconomic status operates as a fundamental cause of physical health disparities, embodying substitutable resources — income, wealth, social connections — that protect health through a shifting set of proximate mechanisms \citep{link_social_1995, phelan_social_2010}. The US literature on the urban-rural dimension of this relationship is characterised in \citet{blumenthal_effects_2002}; the pattern is corroborated in UK data through the Marmot Review evidence base. \\
\addlinespace
Economic\_Capital $\to$ Housing\_Characteristics & Household income and wealth determine the quality of accessible housing stock. Using Family Resources Survey and English Housing Survey data, \citet{waters_housing_2023} show that households in the lowest UK income quintile are substantially more likely to live in homes that are hazardous, difficult to heat, or in poor repair, and that the shrinking pool of properties affordable to housing benefit recipients is of systematically lower quality than the UK housing stock average. \citet{garrison_housing_2025} document similar stratification in the US context. \\
\addlinespace
Economic\_Capital $\to$ Tenure & Wealth and income determine access to owner-occupation in the UK housing market. Social renters in the UK have an after-housing-cost poverty rate of 46\%, private renters 34\%, and mortgagors 11\%, with outright owners at 15\% \citep{waters_housing_2023}. \\
\addlinespace
Economic\_Capital $\to$ Overcrowding & Lower household income constrains consumption of housing space relative to household size. More than two million children from 746,000 families in England live in overcrowded conditions, overwhelmingly concentrated in low-income households \citep{national_housing_federation_overcrowding_2023}. \\
\addlinespace
Economic\_Capital $\to$ Policy\_Scheme\_Access & Means-tested UK energy assistance programmes — most notably Warm Home Discount Core Group 2 — condition eligibility on receipt of qualifying income-related benefits, creating a direct link between area-level economic capital and the reach of energy policy instruments \citep{department_for_energy_security_and_net_zero_warm_2025, department_for_energy_security_and_net_zero_annual_2024}. Expert review proposed reversing this edge (\textit{Policy\_Scheme\_Access $\to$ Economic\_Capital}, on the basis that scheme receipt relieves household budgets); the \textit{Economic\_Capital $\to$ Policy\_Scheme\_Access} orientation is retained because eligibility is conditioned on pre-existing income, but the orientation is flagged for an orientation-sensitivity analysis in which the edge is reversed to confirm that the target estimand is robust to this directional choice. \\
\addlinespace
Economic\_Capital $\to$ Neighbourhood & Household economic resources determine residential neighbourhood through housing market sorting. UK evidence on ethnic and income-based neighbourhood selection demonstrates that higher economic resources facilitate movement to more advantaged neighbourhoods, while residents of constrained means experience limited locational choice \citep{coulter_ethnic_2019, bailey_remaking_2017}. \\
\addlinespace
Economic\_Capital $\to$ Social\_Capital & Economic resources facilitate civic participation through time, financial capacity, and psychosocial stability. A UK analysis of social capital gradients finds a participatory gap between lower- and higher-income groups, though with substantial residual heterogeneity — deprived communities can and do maintain strong social capital through alternative mechanisms, including kinship networks and informal mutual aid \citep{nazroo_civic_2023}. A meta-analysis of socioeconomic gradients in civic engagement corroborates this pattern across high-income countries \citep{neumayr_how_2023}. \\
\addlinespace
Economic\_Capital $\to$ Digital\_Capital & Income determines broadband access, device ownership, and digital literacy. Ofcom's UK digital exclusion evidence finds a strong income gradient in internet use, with approximately one in five adults in households earning under £25{,}000 never using the internet \citep{ofcom_digital_2022}. Broader evidence from European and US studies confirms the socioeconomic stratification of digital access and the emergence of ``relative digital deprivation'' as a secondary inequality \citep{seifert_digital_2024}. \\
\addlinespace
Economic\_Capital $\to$ Substance\_Abuse & Economic deprivation increases the risk of problem substance use. Alcohol-specific death rates in the most deprived decile of UK local authorities are approximately 4.2 times higher than in the least deprived decile \citep{advisory_council_on_the_misuse_of_drugs_vulnerability_2018}; the broader socioeconomic patterning of alcohol-related harm is reviewed in \citet{collins_associations_2016}. \\
\end{longtable}

% =====================================================
% Fuel Poverty edges
% =====================================================

\begin{longtable}{@{}p{3.8cm}p{11.2cm}@{}}
\caption{Outgoing and incoming edges for the Fuel Poverty mediator.}
\label{tab:edges-fuel_poverty}\\
\toprule
\textbf{Edge} & \textbf{Rationale and supporting evidence} \\
\midrule
\endfirsthead
\multicolumn{2}{l}{\textit{Table \ref{tab:edges-fuel_poverty} continued from previous page}}\\
\toprule
\textbf{Edge} & \textbf{Rationale and supporting evidence} \\
\midrule
\endhead
\bottomrule
\endfoot
\multicolumn{2}{@{}l}{\textit{Inbound edges to Fuel Poverty}}\\
\addlinespace

HEE $\rightarrow$ FP &
The energy efficiency of the dwelling stock is a definitional determinant of fuel poverty: the Low Income Low Energy Efficiency (LILEE) metric used in England embeds the efficiency rating directly in its threshold, so areas with less efficient stock mechanically exhibit higher modelled prevalence \citep{department_for_energy_security_and_net_zero_annual_2024}. % VERIFY: confirm report states LILEE methodology, not headline statistics only
Beyond the definitional link, the mechanism is that a more efficient dwelling requires less delivered energy to reach an adequate internal temperature, lowering the expenditure needed for thermal comfort; \citet{sovacool_fuel_2015} reports that home energy-efficiency schemes significantly reduce fuel-poverty prevalence. % VERIFY: confirm prevalence-reduction claim and any figures against full text
The edge is oriented HEE $\rightarrow$ FP because dwelling efficiency is logically prior to the affordability outcome it helps determine \citep{roberts_energy_2008}. \\
\addlinespace

EC $\rightarrow$ FP &
Household income is the second definitional input to fuel-poverty: the low-income component is assessed against an income threshold, so area-level economic capital is constitutive of measured fuel poverty independent of dwelling efficiency \citep{department_for_energy_security_and_net_zero_annual_2024}. The mechanism is direct---for a given energy cost, lower income raises the share of expenditure absorbed by energy and the likelihood of falling below the affordability threshold \citep{galvin_juggling_2024}. This edge is required for FP to be correctly specified; omitting it would misrepresent fuel poverty as a function of efficiency alone. EC already enters the graph as an upstream confounder of HEE (EC $\rightarrow$ HEE); the EC $\rightarrow$ FP edge is conceptually distinct, capturing income's role in affordability \emph{after} efficiency is determined, and the two should not be conflated \citep{roberts_energy_2008}. \\
\addlinespace

C $\rightarrow$ FP &
Climatic heating demand is a recognised driver of fuel poverty: colder areas require more delivered energy to maintain an adequate internal temperature, raising the energy cost faced for a given dwelling efficiency and income. This edge parallels the existing C $\rightarrow$ HEE and C $\rightarrow$ IEQ edges, under which climate already operates through heating demand. The spatial patterning of fuel poverty in England is documented by \citet{robinson_decade_2023}. Energy \emph{price} is excluded as a separate determinant because, in the England-and-Wales LSOA cross-section, tariffs are effectively national and do not vary across units, contributing nothing to identification. \\
\addlinespace

\multicolumn{2}{@{}l}{\textit{Outbound edges from Fuel Poverty}}\\
\addlinespace

FP $\rightarrow$ MH &
Fuel poverty is consistently associated with poorer mental health: in a global scoping review, \citet{khavandi_mental_2024} found a detrimental association in 46 of 47 included studies and organised the pathways into economic, social, behavioural, and environmental mechanisms. The association is observational and confounded by income and housing quality, so it does not by itself establish causation; here those common causes (EC, H) are conditioned, and the residual edge carries the affordability mechanism net of them. The proposed mechanism is psychosocial financial strain---anxiety over unaffordable bills, the heat-or-eat trade-off, and the chronic stress of an under-heated home \citep{longhurst_emotions_2019, grey_short-term_2017}. This edge constitutes the economic/affordability arm of HEE's effect on mental health, running parallel to the physical/thermal arm carried by HEE $\rightarrow$ IEQ $\rightarrow$ MH. \\
\addlinespace

FP $\rightarrow$ PH &
Fuel poverty is linked to adverse physical health, particularly cold-related cardiovascular and respiratory morbidity and excess winter mortality \citep{marmot_review_team_health_2011, maidment_impact_2014}. As with FP $\rightarrow$ MH, the raw association is confounded by income and dwelling characteristics, which the adjustment set addresses; the edge represents the effect of unaffordable energy net of those common causes. The mechanism operates through sustained exposure to inadequate indoor temperatures and the diversion of constrained budgets away from food and health-protective expenditure, with energy-poor households forgoing medicine and healthcare to meet bills \citep{shapira_no_2023}. Evidence is not uniform---\citet{poortinga_health_2018} found no clear reduction in emergency hospital admissions following efficiency investment---so the causal claim is made cautiously. Because PH $\rightarrow$ MH already exists, this edge also opens an indirect FP $\rightarrow$ PH $\rightarrow$ MH sub-path. \\
\addlinespace

FP $\rightarrow$ IEQ &
Fuel poverty degrades indoor environmental quality through a behavioural rationing mechanism that is distinct from the physical-fabric pathway carried by HEE $\rightarrow$ IEQ. Households unable to afford adequate warmth ration heating---under-heating rooms, heating intermittently, or restricting heating to a single room---which lowers internal surface temperatures, and, combined with everyday moisture generation, promotes condensation below the dew point, persistent damp, and mould growth. \citet{sharpe_fuel_2015} find that fuel poverty increases the risk of mould contamination in social housing independently of occupant risk perception and ventilation behaviour, and \citet{liddell_living_2015} document the broader cold-housing pathway. The edge is oriented FP $\rightarrow$ IEQ because affordability-driven heating behaviour is logically prior to the indoor conditions it produces. It is conditionally distinct from the dwelling-form channel (\textit{Housing\_Characteristics $\rightarrow$ IEQ}), the occupancy-load channel (\textit{Overcrowding $\rightarrow$ IEQ}), and the climate heating-demand channel (\textit{Climate $\rightarrow$ IEQ}), and it provides a second route---alongside the existing FP $\rightarrow$ MH and FP $\rightarrow$ PH edges---through which affordability reaches health, namely FP $\rightarrow$ IEQ $\rightarrow$ \{MH, PH\}. \\
\addlinespace

\end{longtable}

% =====================================================
% Urbanicity edges
% =====================================================

\begin{longtable}{@{}p{3.8cm}p{11.2cm}@{}}
\caption{Outgoing edges from Urbanicity.}
\label{tab:edges-urbanicity}\\
\toprule
\textbf{Edge} & \textbf{Rationale and supporting evidence} \\
\midrule
\endfirsthead
\multicolumn{2}{l}{\textit{Table \ref{tab:edges-urbanicity} continued from previous page}}\\
\toprule
\textbf{Edge} & \textbf{Rationale and supporting evidence} \\
\midrule
\endhead
\bottomrule
\endfoot
Urbanicity $\to$ Neighbourhood & Urban–rural classification shapes neighbourhood density, land use, infrastructure provision, and the distribution of local amenities and disamenities. Although many wealthy urban neighbourhoods exist and rural areas exhibit substantial socioeconomic heterogeneity, the urban-rural dimension operates as a structural determinant of neighbourhood composition that sits upstream of deprivation rather than as a simple synonym for it. Evidence of the persistence of area-level urban and deprivation structures in England and Wales over five decades ($r = 0.706$ between 1971 and 2021 deprivation) documents this structural role \citep{lloyd_50-year_2024}. \\
\addlinespace
Urbanicity $\to$ Policy\_Scheme\_Access & Energy policy instruments and household uptake of energy assistance schemes are differentially distributed across urban and rural areas. Rural fuel poverty exhibits distinct characteristics — notably concentration among households off the gas grid and in pre-1919 solid-walled properties — which both shape eligibility for specific retrofit programmes and create barriers to uptake in remote areas \citep{roberts_fuel_2015}. \\
\addlinespace
Urbanicity $\to$ Isolation & Urbanicity affects the geographic and logistical substrate of social isolation, distinct from (though correlated with) the relational pathways captured by \textit{Social\_Capital}. Rural residents face greater physical distances from family, services, and community venues, while urban residents — despite higher population density — report elevated loneliness attributable to weaker neighbourly ties and transient residency patterns \citep{henningsmith_differences_2019, long_loneliness_2024, hussain_loneliness_2023}. The distinction from \textit{Social\_Capital} is that isolation here refers to individual-level geographic and physical access to social contact, whereas social capital captures area-level civic participation and associational density; the two are empirically correlated but represent distinct theoretical constructs. \\
\addlinespace
Urbanicity $\to$ Economic\_Capital & The UK urban–rural pattern in economic capital is heterogeneous and depends on the specific rural sub-category and the measure of income used. Analysis of British Household Panel Survey data finds that urban wages are significantly lower than accessible rural wages but significantly higher than remote rural wages, even after adjusting for observed worker characteristics \citep{gilbert_low_2003}. Urban wage gains are further partially offset by higher urban housing costs and cost-of-living differentials, particularly in London and the South East, such that the net effect on household disposable income after housing costs is ambiguous and area-specific \citep{waters_housing_2023}. The edge is retained to reflect the systematic \textit{structural} relationship between urbanicity and labour market access, while acknowledging that its net magnitude in real terms varies substantially across the urban–rural classification. \\
\addlinespace
Urbanicity $\to$ Education & Socioeconomic disadvantage interacts with rurality in determining educational outcomes. On average, rural pupils slightly outperform urban peers at GCSE in England, but this masks underperformance among rural pupils from deprived backgrounds: once income deprivation is controlled for, rural pupils from the most deprived areas underperform their urban counterparts by approximately 2 percentage points on the 5EM GCSE measure on average, rising to 4–8 percentage points at the extreme ends of deprivation. This pattern is attributed to transport barriers, school sparsity, and narrower post-16 educational choice in rural areas \citep{graham_grass_2024, zahl-thanem_spatial_2024}. \\
\addlinespace
Urbanicity $\to$ Social\_Capital & Rural areas exhibit systematically higher levels of bonding social capital — trust, associational density, and neighbourly contact — than urban areas in England. ONS social capital statistics report that approximately 82\% of rural residents talk regularly to their neighbours compared with 72\% of urban residents \citep{office_for_national_statistics_social_2022}. The mechanism operates through population stability, smaller social units, and the functional necessity of mutual assistance in low-density contexts. \\
\addlinespace
Urbanicity $\to$ Digital\_Capital & A pronounced urban–rural digital divide persists in the UK: over 53\% of deep-rural areas cannot achieve adequate broadband speeds, compared with approximately 5\% in urban areas. This gap reflects the geographic economics of broadband infrastructure deployment and is a structural feature that direct investment has only partially addressed \citep{philip_digital_2017}. \\
\end{longtable}

% =====================================================
% U_AreaHistory (latent) edges
% =====================================================

\begin{longtable}{@{}p{3.8cm}p{11.2cm}@{}}
\caption{Outgoing edges from U\_AreaHistory (latent).}
\label{tab:edges-areahistory}\\
\toprule
\textbf{Edge} & \textbf{Rationale and supporting evidence} \\
\midrule
\endfirsthead
\multicolumn{2}{l}{\textit{Table \ref{tab:edges-areahistory} continued from previous page}}\\
\toprule
\textbf{Edge} & \textbf{Rationale and supporting evidence} \\
\midrule
\endhead
\bottomrule
\endfoot
U\_AreaHistory $\to$ Neighbourhood & Historical planning decisions — slum clearance, council estate construction, industrial siting, motorway routing — physically produced the spatial structure within which present-day neighbourhoods are situated. These decisions shape the housing stock, transport infrastructure, and land-use mix that define neighbourhood boundaries and conditions today, even where present-day deprivation levels do not directly reflect historical patterns. The degree of persistence in area-level structures is substantial: deprivation rankings of areas in England and Wales in 1971 correlate at $r = 0.706$ with deprivation rankings in 2021, demonstrating remarkable stability of relative spatial outcomes over half a century \citep{lloyd_50-year_2024, nardone_relationship_2022}. \\
\addlinespace
U\_AreaHistory $\to$ Economic\_Capital & The legacy of twentieth-century deindustrialisation — colliery closures, steel and manufacturing decline, dock closures — produced persistent intergenerational economic deprivation in affected communities. Causal-inference analysis of UK colliery closures finds that children exposed to pit closures during formative years experienced systematically worse labour-market outcomes throughout adult life, with effects persisting into subsequent generations \citep{brey_death_2024}. The broader structural economic disadvantage of former industrial ``left-behind'' areas is documented by \citet{beatty_beyond_2021}. \\
\addlinespace
U\_AreaHistory $\to$ Education & Historical underinvestment in schools serving deprived areas perpetuates educational disadvantage across generations. UK cohort analyses show that between 1981 and the late 1990s, university graduation rates among young people from the poorest 20\% of families rose by only 3 percentage points, compared to a 26 percentage point rise for those from the richest 20\% — despite substantial overall expansion of higher education during the same period \citep{blanden_educational_2013}. Evidence of the persistence of local-authority-level spending inequalities in English schools into the 2020s is provided in \citet{akanni_inequalities_2025}. \\
\addlinespace
U\_AreaHistory $\to$ PH & Former industrial and deindustrialised ``left-behind'' areas exhibit systematically worse health outcomes, including substantially reduced life expectancy. In the UK's left-behind neighbourhoods, men live on average 3.7 years fewer and women 3.0 years fewer than national averages \citep{case_health_2024}. The US-based redlining literature provides complementary evidence that historical area-level discrimination exerts long-run health effects operating partly through material deprivation and partly through unmeasured structural channels \citep{kraus_historic_2023}. \\
\addlinespace
U\_AreaHistory $\to$ Tenure & The Right to Buy scheme, introduced by the Housing Act 1980, gave eligible council tenants in England and Wales the statutory right to purchase their council home at substantial discount. The policy reshaped UK tenure composition at area level: England's social housing stock fell from approximately 5.5 million units in 1980 to 4.1 million in 2021, with council housing specifically experiencing the largest decline. Cumulative sales over 1980--2023 exceed two million council homes, with sales rates varying sharply across regions and over time in response to discount policy changes \citep{jones_right_2006, disney_right_2017, ministry_of_housing_communities_and_local_government_social_2024}. The tenure landscape produced by these sales — a reduced council sector, expanded owner-occupation, and latterly a substantial private-rental sector composed of former council homes — reflects a historical policy decision rather than a contemporary market outcome. \\
\addlinespace
U\_AreaHistory $\to$ Social\_Capital & Post-war slum clearance and population relocation disrupted established kinship and community networks in UK cities, with documented effects on the social fabric of receiving neighbourhoods and on the capacity for collective action in displaced communities \citep{young_family_1957, crawford_slum_2019}. Subsequent deindustrialisation further fragmented community bonds in former mining and manufacturing areas, where occupational solidarity had served as a key substrate for civic participation. \\
\end{longtable}

% =====================================================
% Neighbourhood outgoing edges
% =====================================================

\begin{longtable}{@{}p{3.8cm}p{11.2cm}@{}}
\caption{Outgoing edges from Neighbourhood (scale-bridging construct).}
\label{tab:edges-neighbourhood}\\
\toprule
\textbf{Edge} & \textbf{Rationale and supporting evidence} \\
\midrule
\endfirsthead
\multicolumn{2}{l}{\textit{Table \ref{tab:edges-neighbourhood} continued from previous page}}\\
\toprule
\textbf{Edge} & \textbf{Rationale and supporting evidence} \\
\midrule
\endhead
\bottomrule
\endfoot
Neighbourhood $\to$ Air\_Pollution & Ambient air pollution varies systematically across areas of differing deprivation in England. Recent LSOA-level analysis covering 2003--2023 finds that the most deprived 20\% of English LSOAs experienced 8\% higher average PM$_{2.5}$ concentrations than the least deprived quintile in 2023, a gap that has widened since 2017 \citep{gadenne_exposure_2024}. Sectoral-emissions analysis at LSOA level finds that all 24 minoritised ethnic groups studied face higher local NOx and PM$_{2.5}$ emissions than socioeconomically matched majority-group populations \citep{gray_evaluating_2024}, consistent with longer-standing evidence from \citet{fecht_associations_2015} and \citet{hajat_socioeconomic_2015}. These findings demonstrate robust spatial correlation between area-level disadvantage and pollution exposure, but do not on their own establish a causal pathway at the area level. A causal interpretation of the scale-bridging edge nonetheless remains well-supported by quasi-experimental evaluations of UK traffic and low-emission interventions: the 2019 London Ultra Low Emission Zone produced a 19.6\% reduction in NO$_2$ at traffic sites and 8.2\% at urban background sites using augmented synthetic control identification \citep{tong_further_2025}; the Birmingham Clean Air Zone achieved 3.4--6.6\% NO$_2$ reductions at roadside sites using the same methodology \citep{liu_assessing_2023}; and evaluations of London Low Traffic Neighbourhoods using difference-in-differences identification found 5.7\% NO$_2$ reductions within intervention areas \citep{yang_evaluation_2022}. These quasi-experimental findings establish that area-level traffic and planning decisions causally determine local pollution exposure, providing mechanism-level support for the scale-bridging claim that upstream factors aggregate through Neighbourhood to shape Air\_Pollution. \\
\addlinespace
Neighbourhood $\to$ Green\_Space & Green and blue space provision is a feature of the local area determined by planning decisions, land-use patterns, and historical investment in public amenities --- all of which aggregate upstream factors into an area-level environmental exposure. A recent UK study of 6{,}791 English MSOAs finds that grassland is inequitably distributed across urban neighbourhoods by Index of Multiple Deprivation: among the most deprived quintile of urban MSOAs in England, a 1\% increase in grassland area is associated with a 37\% reduction in annual preventable deaths (incidence rate ratio 0.63, 95\% CI 0.52--0.76), with the authors explicitly cautioning that this is an observational estimate subject to ecological-inference concerns rather than a causal effect \citep{ngan_inequality_2024}. International evidence of unequal green space distribution is reviewed in \citet{astell-burt_low-income_2014} and \citet{hoffimann_socioeconomic_2017}. \\
\addlinespace
Neighbourhood $\to$ Accessibility & Physical access to healthcare, employment, and services is determined by the spatial distribution of transport networks and service locations at the area level, not by individual household characteristics. The scale-bridging role captures how upstream factors aggregate into area-level accessibility provision. Evidence of the resulting access burden is documented in UK transport research: low-income households in poorly connected areas spend approximately 25\% of their disposable income on commuting costs, compared to an average of around 13\%, reflecting constrained mobility choices that follow from area-level transport provision rather than individual household preference \citep{lucas_inequalities_2019}. Complementary spatial-accessibility evidence is provided by \citet{moreno-monroy_public_2024}. \\
\addlinespace
Neighbourhood $\to$ Neighbourhood\_Stress & Neighbourhood stressors --- crime, disorder, anti-social behaviour, and dereliction --- cluster spatially as area-level exposures. Their distribution reflects the aggregation of upstream economic and historical factors into place-based experiences of unsafety and disorder. In England, 25\% of people in the most deprived decile live in the top 10\% highest-crime areas, compared with 3.1\% of those in the least deprived decile \citep{the_health_foundation_inequalities_2024}. Corroborating evidence on neighbourhood disorder is provided in \citet{polling_influence_2014}. \\
\addlinespace
Neighbourhood $\to$ Housing\_Characteristics & Area-level conditions are associated with systematic differences in the physical housing stock. Overcrowded properties cluster sharply in the most deprived 20\% of LSOAs in England \citep{norman_housing_2022}, and affordable properties in the UK are 15\% more likely to have an EPC rating of D or below than the national housing stock average \citep{waters_housing_2023}. These patterns are cross-sectional and do not on their own establish that area-level aggregation causally produces housing-stock differences. A causal interpretation of the scale-bridging edge rests on the mechanism of ongoing maintenance and investment: the rate at which dwellings are maintained, upgraded, or replaced is shaped by neighbourhood-level rental markets, landlord portfolios, local authority investment, and regeneration decisions that operate at area scale and aggregate upstream household-level investment capacity. Quasi-experimental evidence on this specific mechanism at UK area level is limited, and the edge is retained on the basis of plausibility and theoretical consistency with the scale-bridging framework rather than on direct causal demonstration. \\
\addlinespace
Neighbourhood $\to$ Noise & Exposure to environmental noise --- from traffic, rail, aircraft, and industrial sources --- is an inherently area-level phenomenon determined by the spatial geometry of transport and industrial infrastructure. The scale-bridging role of Neighbourhood captures how upstream land-use and planning decisions aggregate into differential area-level noise exposure. Social inequalities in transport noise are observed across English urban areas, with the most deprived neighbourhoods facing systematically higher exposure \citep{moreno-martos_small-area_2023, tonne_socioeconomic_2018}. \\
\addlinespace
Neighbourhood $\to$ Social\_Capital & Social capital is an intrinsically area-level phenomenon: trust, associational density, and civic participation are properties of a community rather than of an individual. The edge from Neighbourhood captures how upstream factors --- economic, historical, and structural --- aggregate into the area-level context within which social capital forms or fails to form. UK-based nationally representative analysis finds that material deprivation, rather than ethnic or demographic diversity, is the principal neighbourhood-level predictor of reduced social capital \citep{chan_social_2024, forrest_social_2001}, consistent with the aggregation mechanism operating through upstream economic rather than demographic channels. \\
\end{longtable}

% =====================================================
% Neighbourhood_Stress outgoing edges
% =====================================================

\begin{longtable}{@{}p{3.8cm}p{11.2cm}@{}}
\caption{Outgoing edges from Neighbourhood\_Stress (excluding Neighbourhood\_Stress $\to$ MH; see Table~\ref{tab:edges-mh}).}
\label{tab:edges-neighbourhood_stress}\\
\toprule
\textbf{Edge} & \textbf{Rationale and supporting evidence} \\
\midrule
\endfirsthead
\multicolumn{2}{l}{\textit{Table \ref{tab:edges-neighbourhood_stress} continued from previous page}}\\
\toprule
\textbf{Edge} & \textbf{Rationale and supporting evidence} \\
\midrule
\endhead
\bottomrule
\endfoot
Neighbourhood\_Stress $\to$ PH & Chronic exposure to crime, disorder, and perceived unsafety operates on physical health through a sustained physiological stress response. The hypothesised mechanism is allostatic load: repeated activation of the hypothalamic--pituitary--adrenal axis and sympathetic nervous system in response to chronic stressors produces cumulative dysregulation across cardiovascular, metabolic, and immune systems, manifesting as elevated blood pressure, inflammatory markers, and metabolic risk factors. Multi-cohort evidence from the UK Whitehall II, Swiss CoLaus, and Portuguese EPIPorto studies documents that neighbourhood socioeconomic deprivation is associated with higher allostatic load even after adjustment for individual socioeconomic position \citep{ribeiro_neighbourhood_2019}. A systematic review and meta-analysis finds that high allostatic load is associated with a 22\% increase in all-cause mortality and a 31\% increase in cardiovascular-disease mortality \citep{}, linking the neighbourhood-stress-to-allostatic-load pathway to clinically meaningful PH outcomes. Complementary evidence from a systematic review of violence exposure and cardiovascular health identifies consistent associations between exposure to neighbourhood violence and cardiometabolic risk \citep{suglia_violence_2015}. \\
\addlinespace
Neighbourhood\_Stress $\to$ Isolation & Fear of crime produces social withdrawal as an individual adaptive strategy. Residents exposed to crime, perceived unsafety, or violence are more likely to restrict outings, avoid public spaces, and limit interaction with neighbours, reducing opportunities for social contact and increasing loneliness. Evidence from UK longitudinal studies links neighbourhood disorder and perceived unsafety to reduced social interaction and increased loneliness \citep{bowling_how_2007}. \\
\addlinespace
Neighbourhood\_Stress $\to$ Social\_Capital & The mechanism operating through this edge is distinct from the aggregation channel represented by the \textit{Neighbourhood $\to$ Social\_Capital} edge in Table~\ref{tab:edges-neighbourhood}. Whereas that edge captures area-level material deprivation operating on social capital through community stability and civic infrastructure, the present edge operates through a fear-and-withdrawal mechanism: residents exposed to crime and disorder retreat into private space, reduce participation in local associations, and form fewer of the thin trust-based ties with neighbours that constitute social capital. The classic Chicago neighbourhood study of \citet{sampson_neighborhoods_1997} establishes that concentrated disorder undermines collective efficacy independently of material deprivation, and UK-based evidence by \citet{letki_does_2008} reproduces the finding that deprivation-and-disorder --- not ethnic diversity --- is the principal driver of reduced neighbourhood social cohesion. \\
\addlinespace
Neighbourhood\_Stress $\to$ Substance\_Abuse & Neighbourhood stressors increase substance misuse through two complementary channels: a physiological stress pathway in which chronic HPA-axis activation drives self-medication behaviour \citep{sinha_chronic_2008}, and a structural drug-market pathway in which high-crime and disorder-concentrated neighbourhoods exhibit denser local drug availability and lower perceived sanctions against use \citep{felker-kantor_neighbourhood_2019}. Both channels contribute to higher substance-use prevalence in stressed neighbourhoods independently of individual socioeconomic position. \\
\end{longtable}

% =====================================================
% Housing, Climate, and Environmental edges
% =====================================================

\begin{longtable}{@{}p{3.8cm}p{11.2cm}@{}}
\caption{Outgoing edges from Housing\_Characteristics, Climate, Air\_Pollution, Green\_Space, and Noise (environmental mediators, excluding edges in Tables~\ref{tab:edges-hee}--\ref{tab:edges-mh}).}
\label{tab:edges-envmediators}\\
\toprule
\textbf{Edge} & \textbf{Rationale and supporting evidence} \\
\midrule
\endfirsthead
\multicolumn{2}{l}{\textit{Table \ref{tab:edges-envmediators} continued from previous page}}\\
\toprule
\textbf{Edge} & \textbf{Rationale and supporting evidence} \\
\midrule
\endhead
\bottomrule
\endfoot
\multicolumn{2}{@{}l}{\textit{Housing\_Characteristics outgoing edges}}\\
\addlinespace
Housing\_Characteristics $\to$ PH & Poor housing is linked to respiratory disease, cardiovascular stress, and injuries through multiple pathways including structural hazards, inadequate thermal performance, and environmental exposures at the dwelling level. A systematic review identified 19 housing-related risk factors for respiratory disease \citep{wimalasena_housing_2021}, with corroborating evidence from \citet{holden_impact_2023} and \citet{garrison_housing_2025} on the broader cardiovascular and health-injury burden of poor housing. \\
\addlinespace
Housing\_Characteristics $\to$ IEQ & Physical dwelling form --- age, construction era, envelope materials, ventilation design, and maintenance condition --- directly determines indoor temperature regulation, humidity, air quality, and mould risk. Systematic-review evidence identifies a broad range of dwelling-form attributes associated with IEQ outcomes \citep{holden_impact_2023, alfakih_indoor_2025, sharpe_fuel_2015}. \\
\addlinespace
Housing\_Characteristics $\to$ Isolation & High-rise and poorly designed housing reduces opportunities for incidental social interaction and contributes to social withdrawal. A systematic review finds that living in high-rise housing is associated with worse social and mental health outcomes compared with lower-rise residential typologies \citep{barros_social_2019}. \\
\addlinespace
\multicolumn{2}{@{}l}{\textit{Overcrowding outgoing edges (to IEQ)}}\\
\addlinespace
Overcrowding $\to$ IEQ & Residential overcrowding degrades indoor environmental quality through elevated per-occupant moisture generation (respiration, cooking, bathing, laundry) and reduced effective ventilation per occupant, distributed over a fixed building envelope and ventilation capacity. This mechanism is distinct from the dwelling-form pathway captured in the \textit{Housing\_Characteristics} $\to$ \textit{IEQ} edge: the envelope determines ventilation design, while occupancy determines the load placed on it. Empirical evidence links occupant density to elevated indoor humidity, condensation risk, and mould prevalence in UK and international housing \citep{sharpe_fuel_2015, alfakih_indoor_2025, national_housing_federation_overcrowding_2023}. The English Housing Survey operationalises this link through the Housing Health and Safety Rating System (HHSRS) damp-and-mould hazard assessment, which incorporates occupancy density as a contributing factor alongside dwelling condition. \\
\addlinespace
\multicolumn{2}{@{}l}{\textit{Climate outgoing edges}}\\
\addlinespace
Climate $\to$ Housing\_Characteristics & Extreme weather --- flooding, heat, and storms --- degrades the structural integrity and condition of dwellings. In England, approximately 6.3 million properties are currently at flood risk, with projected increases by 2050 as climate change intensifies \citep{dawson_systems_2018, hayles_quantifying_2022}. Coastal and geographic exposure operates through the same channel and is, following expert review, represented within the exogenous \textit{Climate} construct rather than as a separate node: dwellings in coastal and highly exposed LSOAs face elevated coastal-flood and storm-surge risk, wind-driven rain, and accelerated salt- and wind-weathering of the building fabric, which raise rates of fabric deterioration and repair need. \\
\addlinespace
Climate $\to$ IEQ & Extreme outdoor temperatures propagate to indoor environments with substantial time lags and amplification dependent on dwelling fabric. Self-reported evidence from a national online survey of approximately 1{,}580 UK households conducted following the record 2022 heatwave, when outdoor temperatures exceeded $40^\circ$C for the first time in the UK, reports a rise in household-reported indoor overheating from 20\% (2011) to 82\% (2022) \citep{khosravi_nation_2025}. These subjective figures are higher than those obtained from objective indoor-temperature monitoring: using the 2011 Energy Follow-Up Survey with CIBSE TM59 criteria, approximately 26\% of English dwellings exceeded Criterion~2 \citep{pathan_monitoring_2017}, with higher prevalence projected under warming scenarios \citep{symonds_home_2021}. Climate additionally drives indoor damp through moisture ingress: wind-driven (driving) rain penetrates the building envelope---disproportionately in exposed coastal and high-rainfall western areas and in solid-walled stock---raising wall moisture content, condensation risk, and mould, a mechanism recognised in UK building-fabric and exposure guidance \citep{hayles_quantifying_2022}. Flooding additionally introduces damp and mould, with parallel IEQ consequences for indoor air quality \citep{mansouri_impact_2022}. \\
\addlinespace
\multicolumn{2}{@{}l}{\textit{Air\_Pollution outgoing edges}}\\
\addlinespace
Air\_Pollution $\to$ PH & Long-term exposure to ambient air pollution causes respiratory and cardiovascular disease and premature mortality. A multi-country analysis of 854 European cities estimates that each 10\,\textmu g/m$^{3}$ increase in PM$_{2.5}$ is associated with approximately an 8\% increase in all-cause mortality \citep{khomenko_long-term_2024}. \\
\addlinespace
Air\_Pollution $\to$ IEQ & Outdoor particulate pollution infiltrates indoor environments through the building envelope, with the degree of infiltration dependent on airtightness, ventilation, and occupant behaviour. A systematic review of 77 studies covering over 4{,}000 homes (approximately 1{,}000 in the United States and 150 across Europe) reports a mean infiltration factor for PM$_{2.5}$ of 0.55, indicating that on average more than half of ambient PM$_{2.5}$ penetrates and persists indoors \citep{chen_review_2011}. UK-specific infiltration data is limited, though dwelling airtightness and ventilation patterns in the English housing stock suggest comparable infiltration behaviour \citep{shrubsole_impacts_2016}. \\
\addlinespace
\multicolumn{2}{@{}l}{\textit{Green\_Space outgoing edges}}\\
\addlinespace
Green\_Space $\to$ PH & Access to green space reduces cardiovascular-disease risk and all-cause mortality through pathways including physical activity, stress reduction, and heat mitigation. A meta-analysis synthesising evidence from over 100 million individuals reports a 2--3\% reduction in cardiovascular-disease mortality per 0.1-unit increase in Normalised Difference Vegetation Index (NDVI), a standardised remote-sensing measure of vegetation density ranging from $-1$ (non-vegetated surfaces) to $+1$ (dense vegetation) \citep{liu_green_2022, rojas-rueda_green_2019, kruize_urban_2019}. A 0.1-unit increase corresponds, for example, to the difference between low-density suburban vegetation and an area containing established parks or tree cover. \\
\addlinespace
Green\_Space $\to$ IEQ & Urban vegetation moderates the urban heat island effect, lowering ambient air temperatures by approximately 1--7$^\circ$C depending on park size, vegetation type, and climate \citep{bowler_urban_2010, aram_urban_2019, iungman_cooling_2023}. The indirect effect on indoor temperatures operates through reduced outdoor-to-indoor heat transfer and is sensitive to building envelope and ventilation regime. \\
\addlinespace
Green\_Space $\to$ Air\_Pollution & Urban trees and vegetation remove particulate matter through dry deposition and absorb gaseous pollutants through stomatal uptake. US urban trees are estimated to remove approximately 711{,}000 tonnes of pollutants annually \citep{nowak_modeled_2013, maher_impact_2013}. \\
\addlinespace
\multicolumn{2}{@{}l}{\textit{Noise outgoing edges}}\\
\addlinespace
Noise $\to$ PH & Chronic exposure to environmental noise --- principally from road, rail, and aircraft sources --- causes cardiovascular disease through sustained sympathetic activation and disrupted sleep. A WHO-commissioned quantitative review of the environmental-noise evidence base finds consistent dose-response associations between community noise exposure and cardiovascular morbidity \citep{van_kempen_quantitative_2012}, with mechanism-level support from \citet{munzel_transportation_2021}. Evidence on occupational-noise exposure, which is distinct from the environmental-noise regime that characterises LSOA-level exposure, identifies approximately 1.81$\times$ higher hypertension risk at workplace noise levels $\geq 80$\,dB(A) \citep{bolm-audorff_occupational_2020}, providing complementary mechanism-level evidence on the noise--cardiovascular pathway while noting the higher exposures typical of occupational settings. \\
\addlinespace
Noise $\to$ IEQ & External noise penetrates indoor environments primarily through windows, doors, and building envelope weaknesses. UK building standard BS~8233:2014 specifies indoor noise targets to protect occupants from outdoor noise ingress \citep{british_standards_institution_bs_2014, world_health_organization_environmental_2018}. \\
\end{longtable}

% =====================================================
% HEE outgoing, IEQ outgoing, and Social/Psychological edges
% =====================================================

\begin{longtable}{@{}p{3.8cm}p{11.2cm}@{}}
\caption{Groups 12--14 --- HEE outgoing, IEQ outgoing, Social/Psychological Mediator, and Demographics outgoing edges.}
\label{tab:edges-social}\\
\toprule
\textbf{Edge} & \textbf{Rationale and supporting evidence} \\
\midrule
\endfirsthead
\multicolumn{2}{l}{\textit{Table \ref{tab:edges-social} continued from previous page}}\\
\toprule
\textbf{Edge} & \textbf{Rationale and supporting evidence} \\
\midrule
\endhead
\bottomrule
\endfoot
\multicolumn{2}{@{}l}{\textit{HEE outgoing edges}}\\
\addlinespace
HEE $\to$ IEQ
& Energy efficiency improvements alter indoor environments in complex ways. Properly implemented retrofits that incorporate adequate ventilation yield net IEQ gains; however, retrofits undertaken without ventilation provision can trap indoor pollutants and worsen air quality. Empirical evidence documents both the potential health benefits and the IEQ risks depending on whether ventilation is upgraded alongside envelope improvements \citep{hamilton_health_2015, shrubsole_100_2014, fisk_association_2020, poortinga_health_2018}. \\
\addlinespace
HEE $\to$ Noise
& Double- and triple-glazing, solid wall insulation, and improvements in airtightness reduce the transmission of external noise through the building envelope. Standard double glazing achieves a reduction of approximately 30--35~dB relative to single glazing, with further gains from triple glazing and improved frame sealing \citep{tong_full-scale_2015, hongisto_satisfaction_2015}. \\
\addlinespace
\multicolumn{2}{@{}l}{\textit{IEQ outgoing edges}}\\
\addlinespace
IEQ $\to$ PH
& Adverse indoor environmental conditions --- including dampness, mould, cold temperatures, and elevated concentrations of indoor pollutants --- cause respiratory disease, cardiovascular strain, and excess winter mortality. Meta-analyses of controlled and observational evidence report odds ratios of 1.34--1.75 for respiratory outcomes associated with damp and mouldy housing \citep{fisk_meta-analyses_2007, world_health_organization_who_2009, marmot_review_team_health_2011}. \\
\addlinespace
\multicolumn{2}{@{}l}{\textit{Social and Psychological Mediator edges}}\\
\addlinespace
Accessibility $\to$ PH & Physical access to healthcare determines the timeliness and continuity of treatment for acute and chronic physical illness, operating through deferred diagnosis, missed appointments, and discontinuous management of long-term conditions. A scoping review of distance and travel time to health services in rural and remote areas finds that greater distance and travel time are associated with delayed presentation, lower service utilisation, and worse health outcomes, with effects concentrated among elderly and chronically ill populations \citep{mseke_impact_2024}. UK health-service inequalities research documents the same area-level access burden \citep{lowther-payne_understanding_2023}. This pathway is mechanistically distinct from \textit{Accessibility $\to$ MH} (Table~\ref{tab:edges-mh}): the present edge captures the consequences of constrained physical-healthcare access for physical morbidity and mortality, rather than the psychological effects of restricted participation. Because \textit{PH $\to$ MH} already exists, this edge also opens an indirect \textit{Accessibility $\to$ PH $\to$ MH} sub-path running parallel to the direct \textit{Accessibility $\to$ MH} channel. \\
\addlinespace
Social\_Capital $\to$ Isolation
& Strong community networks reduce social isolation by facilitating incidental interaction, reinforcing trust, and providing a sense of belonging. Evidence from cross-sectional and longitudinal studies consistently identifies higher social capital as protective against loneliness and social withdrawal \citep{nyqvist_association_2016, holtlunstad_social_2024}. \\
\addlinespace
Social\_Capital $\to$ Substance\_Abuse
& Community cohesion protects against substance misuse through the enforcement of social norms and the provision of informal support networks. A systematic review of alcohol interventions finds that higher social capital is associated with significantly reduced binge drinking (adjusted OR~0.38), with complementary evidence on the broader protective role of community ties \citep{bryden_systematic_2013, weitzman_risk_2005}. \\
\addlinespace
Digital\_Capital $\to$ Isolation
& Digital exclusion limits access to online social interaction and public services, increasing the risk of loneliness. An analysis of the English Longitudinal Study of Ageing (ELSA) finds that digital exclusion is associated with higher odds of loneliness (OR~1.30), with comparable evidence from additional cohort studies \citep{zhong_digital_2025, yan_association_2024}. \\
\addlinespace
Substance\_Abuse $\to$ PH
& Alcohol and drug misuse cause liver disease, cardiovascular damage, cancer, and premature mortality through direct physiological mechanisms. A Global Burden of Disease analysis estimates that 2.8 million deaths globally were attributable to alcohol consumption in 2016 \citep{gbd_2016_alcohol_collaborators_alcohol_2018}. \\
\addlinespace
\multicolumn{2}{@{}l}{\textit{Demographics outgoing edges}}\\
\addlinespace
Demographics $\to$ Housing
& Ethnicity, age, and household composition influence housing quality through differential access to housing markets and documented discrimination in allocation and lending. Pakistani and Black Other households experience overcrowding rates exceeding 20\%, reflecting structurally unequal outcomes in housing access \citep{resolution_foundation_heritage_2025, brimblecombe_mapping_2024}. \\
\addlinespace
Demographics $\to$ Overcrowding
& Household composition and ethnicity are among the strongest determinants of overcrowding. 2021 Census data show that Black households faced an overcrowding rate of 16.1\%, compared with 2.5\% for White households, illustrating the magnitude of demographic disparities in housing density \citep{office_for_national_statistics_overcrowding_2023, dezateux_inequalities_2025}. \\
\addlinespace
Demographics $\to$ Economic\_Capital
& Gender, ethnicity, and age systematically shape earnings through labour market discrimination and structural barriers to employment and progression. Bangladeshi employees face the largest documented negative ethnicity pay gap ($-$16\%), with parallel penalties associated with gender and age \citep{briggs_understanding_2020, office_for_national_statistics_ethnicity_2023, goldin_understanding_1992}. \\
\addlinespace
Demographics $\to$ Education
& Socioeconomic background, gender, and ethnicity are significant predictors of educational attainment. At age 16, the socioeconomic achievement gap is approximately twice the size of the largest ethnic attainment gap and six times the gender gap, underscoring the dominance of socioeconomic origin in shaping educational outcomes \citep{strand_ethnicity_2014, li_entrenched_2021}. \\
\end{longtable}

% =====================================================
% Excluded edges
% =====================================================

\begin{longtable}{@{}p{3.8cm}p{11.2cm}@{}}
\caption{Edges considered but excluded from the DAG, with the prima facie case for inclusion and the reason for omission.}
\label{tab:excluded-edges}\\
\toprule
\textbf{Excluded edge} & \textbf{Prima facie case and reason for exclusion} \\
\midrule
\endfirsthead
\multicolumn{2}{l}{\textit{Table \ref{tab:excluded-edges} continued from previous page}}\\
\toprule
\textbf{Excluded edge} & \textbf{Prima facie case and reason for exclusion} \\
\midrule
\endhead
\bottomrule
\endfoot
Education $\to$ HEE & Higher educational attainment might plausibly increase retrofit uptake through improved energy literacy, awareness of subsidy programmes, and capacity to evaluate long-term returns on efficiency investments. However, empirical analysis of British household survey data finds that a BA-degree indicator is not a statistically significant predictor of retrofit investment once income and dwelling characteristics are controlled for \citep{trotta_factors_2018}. Cross-national studies reporting positive education-retrofit associations attribute them predominantly to education's correlation with income rather than to an independent information channel. We therefore model the effect of education on HEE as operating entirely through \textit{Education $\to$ Economic\_Capital $\to$ HEE}, and omit the direct edge. \\
\addlinespace
Demographics $\to$ HEE & Fuel poverty incidence and EPC distribution differ systematically by ethnicity, age, and household composition in England, making a direct demographic effect on HEE intuitive. Direct empirical tests of this question, however, consistently find that demographic variables' explanatory power over HEE operates predominantly through economic and housing mediators already represented in the DAG. A recent machine learning analysis at LSOA level in England and Wales using 2021 Census and EPC data — the scale and setting of the present analysis — finds that socio-demographic variables collectively explain approximately one-third of the variance in HEE, but that the dominant predictors are the National Statistics Socio-Economic Classification, household deprivation, and median age, while ethnicity and religious affiliation exert comparatively modest influence \citep{au_yong_assessing_2026}. The leading predictors in that analysis are direct proxies for \textit{Economic\_Capital} rather than for \textit{Demographics} as constructed in this DAG, and the analysis does not condition on dwelling characteristics, which separate work has consistently found to be the dominant determinant of HEE \citep{trotta_factors_2018, trotta_determinants_2018, sheridan_energy_2023, buyuklieva_variations_2024}. These findings are consistent with earlier evidence that socio-demographic characteristics are substantially less important than dwelling-related characteristics in predicting retrofit uptake in English households once income is controlled for. Demographics' effect on HEE is therefore adequately captured through the established indirect pathways \textit{Demographics $\to$ Housing\_Characteristics $\to$ HEE}, \textit{Demographics $\to$ Economic\_Capital $\to$ HEE}, and \textit{Demographics $\to$ Tenure $\to$ HEE}, and a direct edge would risk double-counting effects already transmitted through these mediators. Fuel poverty, a joint function of income and HEE \citep{department_for_energy_security_and_net_zero_annual_2024}, is not equivalent to HEE and is treated accordingly. \\
\addlinespace
HEE $\to$ Economic\_Capital & Energy efficiency improvements reduce household energy bills, which at the household level represents an increase in disposable income. Including this edge would make HEE a cause of downstream economic capital and would substantially complicate identification. At the LSOA area level, however, the bill-saving effect of HEE is already represented within the DAG through the affordability channel to mental health (\textit{HEE $\to$ IEQ} captures thermal comfort; the financial stress relief operating through reduced fuel poverty is captured in the HEE $\to$ MH mechanism documented in \citet{willand_towards_2015} and \citet{grey_short-term_2017}). Adding a direct \textit{HEE $\to$ Economic\_Capital} edge would conflate the exposure of interest with one of its downstream consequences, rendering the target estimand ambiguous. The edge is omitted, and the affordability mechanism is acknowledged as part of the HEE--MH pathway rather than as a modifier of the confounder set. \\
\addlinespace
Housing\_Characteristics $\to$ MH & A direct effect of physical dwelling form on mental health is suggested by prospective evidence that persistent poor housing conditions predict psychological distress over multiyear horizons \citep{pevalin_impact_2017, singh_housing_2019}. In the present DAG, however, \textit{Housing\_Characteristics} is defined as physical dwelling form (age, construction era, type, materials), while housing \textit{quality} — damp, mould, cold, and indoor pollutants, which are the proximal drivers of the mental health effects these studies document — is represented by the \textit{IEQ} node. The mental health consequences of physical housing are therefore transmitted through \textit{Housing\_Characteristics $\to$ IEQ $\to$ MH}, with additional pathways via \textit{Overcrowding} and \textit{Isolation}. A direct edge would double-count the IEQ-mediated effect that constitutes the dominant evidence base in the housing-mental-health literature. \\
\addlinespace
Economic\_Capital $\to$ Education & Parental economic resources shape children's educational opportunities through school quality, tutoring, and sustained enrolment; this intergenerational mechanism is well-documented \citep{blanden_educational_2013}. Including this edge would, however, introduce a cycle with the existing \textit{Education $\to$ Economic\_Capital} edge, which captures the adult labour-market returns to schooling. DAG formalism requires a directional choice. The \textit{Education $\to$ Economic\_Capital} direction is retained because (i) at any single cross-section at LSOA level, the dominant data-generating process for working-age earnings is the labour-market-returns mechanism; (ii) the intergenerational reverse-direction mechanism operates on timescales outside the temporal scope of the present analysis; and (iii) the latent node \textit{U\_AreaHistory} already absorbs part of the area-level intergenerational persistence that would otherwise require a \textit{Economic\_Capital $\to$ Education} edge \citep{lloyd_50-year_2024}. The bidirectional nature of the underlying relationship is an acknowledged limitation rather than a modelling error. \\
\addlinespace
MH $\to$ PH & Common mental disorders increase physical morbidity through neuroendocrine, inflammatory, and behavioural pathways; mediation analyses using UK data find significant effects in both the physical-to-mental and mental-to-physical directions \citep{ohrnberger_relationship_2017, berk_comorbidity_2023}. The relationship is genuinely bidirectional, but the DAG framework requires a single directional choice. We retain \textit{PH $\to$ MH} and omit \textit{MH $\to$ PH} because, for the target estimand (the effect of HEE on MH), the \textit{PH $\to$ MH} direction is what creates backdoor paths through physical health and therefore what matters for confounder identification. Prioritising the direction that threatens identification is the more conservative modelling choice; the reverse pathway, while real, would introduce a cycle without altering the adjustment set for our exposure-outcome pair. The implicit feedback loop is an acknowledged limitation of a cross-sectional DAG applied to a substantively dynamic system. \\
\addlinespace
Neighbourhood\_Stress $\to$ Green\_Space & Fear of crime and neighbourhood disorder reduce residents' behavioural engagement with local parks and green infrastructure, particularly among women and ethnic minorities \citep{sreetheran_socio-ecological_2014, branas_citywide_2018}. In the present DAG, however, \textit{Green\_Space} is defined as the area-level provision and quality of parks, vegetation, and blue space, measured at LSOA level through land-cover and density indicators rather than use or engagement data. The behavioural effect of neighbourhood stress on green-space engagement is conceptually distinct from provision, and its mental health consequences are captured through two existing pathways: \textit{Neighbourhood\_Stress $\to$ Isolation $\to$ MH} (reduced social interaction) and the direct \textit{Neighbourhood\_Stress $\to$ MH} channel (chronic stress). A direct edge to \textit{Green\_Space} would require redefining that node to include use, breaking consistency with Tables~\ref{tab:edges-mh} and~\ref{tab:edges-neighbourhood}. \\
\addlinespace
Climate $\to$ HEE & Regional climate plausibly shapes heat-demand profiles and thus the marginal value of efficiency measures, and longitudinal monitoring of UK social housing shows spatial variation in heating demand attributable to regional weather \citep{aragon_influence_2022}. In the present DAG, however, \textit{HEE} is the installed physical efficiency of the dwelling stock (insulation, glazing, heating-system quality, airtightness), which in an England-and-Wales LSOA cross-section is determined by dwelling fabric, household income, tenure, and policy uptake rather than by regional climate. Climate operates on heating \textit{demand} and on the indoor and affordability consequences of a given efficiency level — channels already represented by \textit{Climate $\to$ IEQ} and \textit{Climate $\to$ Fuel\_Poverty}. A direct \textit{Climate $\to$ HEE} edge would conflate heating demand with dwelling efficiency and risk double-counting the demand-side pathway; it is therefore omitted, following expert review. The summertime-overheating interaction between climate and the building envelope is retained through \textit{Climate $\to$ IEQ}. \\
\end{longtable}